\documentclass[11pt,a4paper]{article}
\usepackage[utf8]{inputenc}
\usepackage[T1]{fontenc}
\usepackage{amsmath, amssymb, amsfonts}
\usepackage{graphicx}
\usepackage{geometry}
\geometry{a4paper, margin=1in}
\usepackage{hyperref}
\usepackage{xcolor}
\usepackage{listings}

\title{\textbf{Understanding smolVLA: Training and Evaluation Details}}
\author{Finetuning, Losses, and the Overfitting Paradigm}
\date{\vspace{-5ex}}

\begin{document}

\maketitle

\section{Introduction}
Training a \textbf{smolVLA} model involves fine-tuning a pretrained Vision-Language-Action model on a specific robot dataset. This process differs from standard computer vision or NLP training in how it measures success and how it handles data validation. 

This document clarifies the interpretation of training metrics (the Loss) and the specific evaluation strategy used in the LeRobot framework.

\section{Interpreting the Training Loss}
During training, the primary metric displayed in the logs (e.g., via \texttt{wandb} or terminal output) is the \textbf{Training Loss}. As established in the Flow Matching framework, this loss is the \textbf{Mean Squared Error (MSE)} between the target velocity vector $u_\tau$ and the model's predicted velocity $v_\tau$.

\subsection{What does a loss value mean?}
Let's analyze a typical loss value, such as $\mathcal{L} = 0.02$.

\begin{enumerate}
    \item \textbf{The Mathematical Definition:} The loss is the average of squared differences across all elements of the action chunk. For a standard smolVLA configuration ($50$ timesteps $\times 32$ joints), each loss value is an average over $1,600$ individual predictions.
    \[ \text{MSE} = \frac{1}{N} \sum (v_{i,j} - u_{i,j})^2 \]
    
    \item \textbf{Root Mean Squared Error (RMSE):} To understand the "average error" in the same units as the velocity, we take the square root of the loss:
    \[ \text{RMSE} = \sqrt{0.02} \approx 0.1414 \]
    
    \item \textbf{Scale Comparison:} 
    \begin{itemize}
        \item The target velocity $u_\tau$ is defined as $\text{Noise} - \text{Action}$. 
        \item Since the \textbf{Noise} is sampled from a standard Gaussian $\mathcal{N}(0, 1)$, its values typically range from $-3$ to $3$.
        \item The \textbf{Actions} are normalized by the LeRobot dataset statistics, typically mapping them to the $[-1, 1]$ range.
        \item Therefore, the target $u_\tau$ often spans a range of approximately \textbf{$-4$ to $4$}.
    \end{itemize}
    
    \item \textbf{Intuition:} If your loss is $0.02$, it means that on average, the model's prediction for each of the $1,600$ velocity components is off by only \textbf{$\approx 0.14$}. Given that the target values can be as large as $3$ or $4$, an error of $0.14$ represents a very high degree of precision in predicting the flow field.
\end{enumerate}

\section{The Question of Overfitting and Validation}
In standard machine learning, we carefully split data into \textit{Training}, \textit{Validation}, and \textit{Test} sets to detect overfitting. However, the \texttt{lerobot-train} pipeline for offline imitation learning often behaves differently.

\subsection{Is there an Offline Validation Set?}
By default, the \texttt{lerobot\_train.py} script \textbf{does not perform an offline train/test split} of the demonstration data to compute a "Validation Loss." 

The training loop cycles through the entire provided dataset continuously. There are several reasons for this design choice:
\begin{enumerate}
    \item \textbf{Data Scarcity:} Robot demonstration datasets are often small (50 to 100 episodes). Splitting off 20\% of this precious expert data for a validation loss would significantly reduce the diversity of behaviors the model can learn.
    \item \textbf{Loss vs. Success:} In Robotics, \textbf{Offline Loss is a poor proxy for Online Success.} A model can have a very low validation loss (it predicts the expert's next move well on paper) but fail miserably in the real world because it hasn't learned to recover from small deviations (the "compounding error" problem).
\end{enumerate}

\subsection{If there's no Val Set, how do we evaluate?}
Evaluation in LeRobot is performed \textbf{"In-the-Loop"} rather than \textbf{"On-Paper"}:
\begin{itemize}
    \item \textbf{Environmental Evaluation (\texttt{eval\_freq}):} Every $N$ steps (e.g., 20,000), the trainer pauses and runs the current policy in a \textbf{Simulator} (e.g., Gym, Mujoco, or Isaac Gym).
    \item \textbf{True Metrics:} Instead of looking at MSE, we look at the \textbf{Success Rate} (percentage of tasks completed) and \textbf{Average Reward}. This is the only "validation" that truly matters for a robot.
    \item \textbf{Real-World Limitation:} For real-world datasets where no simulator exists, there is indeed no automatic validation. Practitioners usually train for a fixed number of steps, then manually test the checkpoints on the physical robot.
\end{itemize}

\subsection{Is Overfitting a Problem?}
In Imitation Learning (Behavioral Cloning), \textbf{overfitting is often precisely the goal.} We want the model to mimic the expert as perfectly as possible. 

The "Flow Matching" framework naturally mitigates the traditional dangers of overfitting:
\begin{itemize}
    \item By sampling new random noise matrices for every single batch, we are effectively performing massive \textbf{Data Augmentation}. The model never sees the exact same "noisy chunk" twice, which forces it to learn the underlying \textbf{Vector Field} rather than memorizing specific sequences.
    \item Even if the model "overfits" the training demonstrations, its ability to denoise starting from \textbf{Random Noise} at inference time provides it with inherent robustness and stochastic variety.
\end{itemize}

\section{Action Scaling: From Degrees to Latents}
A natural question arises: if robot sensors report angles from $-180^\circ$ to $180^\circ$, how do we add noise to those large numbers?

The answer is that the actions are \textbf{Normalized} before the training loop ever sees them.
\begin{enumerate}
    \item \textbf{Statistics Collection:} Before training starts, LeRobot calculates the \texttt{mean} and \texttt{std} (standard deviation) for every single joint and motor across the entire dataset.
    \item \textbf{The Transform:} During training, the raw actions are transformed:
    \[ \text{Action}_{norm} = \frac{\text{Action}_{raw} - \text{mean}}{\text{std}} \]
    \item \textbf{Stable Noising:} This maps the actions to a range roughly around $[-1, 1]$. Now, when we add Gaussian noise ($\mathcal{N}(0, 1)$), the noise is of a similar magnitude to the signal. This ensures that the model can actually "see" the noise it needs to remove.
\end{enumerate}

\section{The Simulator: Beyond the MSE}
You asked if the simulator just measures MSE on noised chunks. The answer is \textbf{No}—that would be "Offline Evaluation." The simulator evaluation is much more powerful: it is a \textbf{Rollout}.

\subsection{How the Simulator Rollout Works}
When \texttt{lerobot-train} pauses for evaluation, it doesn't just look at expert data. It lets the model "drive" the robot in a virtual environment:
\begin{enumerate}
    \item \textbf{Reset:} The simulator (e.g., Mujoco or Gym) resets to a fresh starting state.
    \item \textbf{The Selection Loop:} 
    \begin{itemize}
        \item The model receives a \textit{live} image and state from the simulator.
        \item The model generates a new action chunk using the full iterative Flow Matching process (10 steps of Euler integration).
        \item The \textit{first action} of that chunk is applied to the simulator's motors.
        \item The simulator physics engine calculates the new state (the objects move, the arm stays or falls).
    \end{itemize}
    \item \textbf{Completion:} This continues until the episode ends or the robot times out.
\end{enumerate}

\subsection{Success vs. Loss}
The simulator measures two key metrics that MSE cannot capture:
\begin{itemize}
    \item \textbf{Success Rate:} It checks the actual outcome (e.g., "Is the block inside the red bowl?").
    \item \textbf{Compounding Errors:} In a rollout, if the model makes a small mistake at Step 10, the robot ends up in a position it never saw in the training data. The rollout tests if the model is \textbf{robust} enough to recover and still finish the task.
\end{itemize}

\subsection{Optional Evaluation: Skipping the Simulator}
It is important to note that the Simulator Evaluation is \textbf{optional}. If the configuration setting \texttt{env} is set to \texttt{null} (as it was in the \texttt{smolvla\_red\_block\_finetune} run), the trainer will simply skip the rollout phase.

In such cases, the user relies entirely on the \textbf{Training Loss} to monitor progress. This is common when:
\begin{itemize}
    \item No matching virtual simulator exists for the specific physical setup.
    \item The user wants to speed up training by avoiding the overhead of running episodes.
    \item The goal is to first ensure the model can converge on the expert manifold before testing live behaviors.
\end{itemize}

\section{Summary}
\begin{itemize}
    \item \textbf{Normalization:} We scale the $-180/180$ range down to $\approx [-1, 1]$ before adding noise so the math stays stable.
    \item \textbf{Simulator:} It's an "Online" test. The robot actually tries to do the task. This is the ultimate "Validation" of the model's intelligence.
    \item \textbf{Flexibility:} If no simulator is available (e.g., \texttt{env: null}), the training continues as a pure offline process, evaluated purely via training loss.
\end{itemize}

\end{document}
